\section{Introdução}
O envelhecimento cutâneo reúne mudanças estruturais e funcionais que se acumulam ao longo da vida. Fatores intrínsecos, ligados à idade e às características individuais, interagem com exposições ambientais, sobretudo a radiação ultravioleta. Essa combinação pode se manifestar em rugas, perda de elasticidade, ressecamento e alterações de pigmentação. Compreender esses processos ajuda a orientar a busca por intervenções que preservem a integridade da pele e atenuem seus efeitos \parencite{liu2023,papaccio2022}.

Parte dessas mudanças envolve a remodelação da matriz extracelular, que fornece suporte à derme e contribui para as propriedades mecânicas da pele. A produção de melanina também participa da pigmentação. Os alvos selecionados representam processos diferentes: a metaloproteinase de matriz 1 (MMP-1) inicia a clivagem dos colágenos fibrilares I e III, importantes para a resistência da derme. MMP-2 degrada gelatina e colágeno IV da membrana basal; MMP-9 também degrada gelatina e componentes da membrana basal, participando da remodelação da interface entre epiderme e derme \parencite{liu2025,papaccio2022}.

A elastase de neutrófilos degrada elastina das fibras elásticas, que ajudam a pele a recuperar sua forma após deformação. A hialuronidase 1 hidrolisa o ácido hialurônico, componente da matriz que retém água e contribui para a hidratação; sua ação participa da renovação dessa molécula. Já a proteína relacionada à tirosinase 1 (TYRP1) atua em melanócitos e está associada à melanogênese e à pigmentação. Em humanos, porém, TYRP1 não apresenta a atividade oxidativa característica da tirosinase. Assim, o painel reúne alvos ligados à estrutura, à hidratação e à pigmentação, sem pressupor que tenham funções equivalentes ou que sua modulação produza o mesmo efeito na pele \parencite{papaccio2022,chao2007,lai2017,liu2025,macdonald2002}.

As microalgas oferecem materiais diversos para a pesquisa de compostos bioativos \parencite{khemiri2023}. Em \textit{Tetraselmis chuii}, foram estudados extratos, biomassa e vesículas extracelulares. Essas vesículas são pequenas partículas liberadas pelas células no meio de cultivo; cada uma é envolvida por uma bicamada lipídica que transporta moléculas. \textcite{adamo2024} denominaram nanoalgossomos as vesículas isoladas de culturas de \textit{T. chuii} e relataram sua biocompatibilidade, além de propriedades antioxidantes e anti-inflamatórias. Elas são distintas das substâncias poliméricas extracelulares (SPE): enquanto as vesículas são partículas delimitadas por membrana, a SPE forma uma matriz extracelular composta principalmente por polissacarídeos \parencite{delattre2016}.

Cokdinleyen et al. investigaram compostos de extratos com potencial neuroprotetor. Em modelos de células humanas da pele, Gargano et al. observaram que os nanoalgossomos reduziram alterações associadas ao estresse oxidativo e à senescência após exposição à radiação ultravioleta B (UVB), além de diminuir o conteúdo de melanina em melanócitos estimulados \parencite{adamo2024,cokdinleyen2024,gargano2026}. 

\Needspace{12\baselineskip}
A SPE forma uma matriz extracelular associada à superfície das células ou liberada no meio de cultivo. É composta principalmente por polissacarídeos, além de proteínas, lipídios e ácidos nucleicos; sua fração de carboidratos reúne diferentes açúcares, e a composição varia conforme a espécie e as condições de cultivo \parencite{delattre2016}. Em exopolissacarídeos de \textit{T. suecica}, foram identificados glicose, galactose, manose, arabinose, ramnose, xilose e fucose, além dos ácidos glucurônico e galacturônico \parencite{parra2020}. Na biomassa de \textit{T. chuii}, glicose, galactose, manose, arabinose e ramnose foram descritas por \citeauthor{brown1991} (\citeyear{brown1991}, \textit{apud} \citeauthor{pereira2019}, \citeyear{pereira2019}).

Com base nas evidências descritas, foram selecionados 17 compostos para triagem, distribuídos em monossacarídeos, derivados de açúcares e fenólicos, entre eles glicose-6-sulfato, galactose-6-sulfato, ácido gálico e ácido ferúlico. A diversidade do conjunto permite examinar associações com proteínas cutâneas e compará-las com ligantes comerciais e, quando disponíveis, coocristalizados. Unidades de glicose sulfatadas foram relatadas em polissacarídeos de \textit{Porphyridium} sp.; os ácidos gálico e ferúlico foram identificados em extratos de biomassa de \textit{T. tetrahele} \parencite{geresh2009,mai2026}.

O docking molecular permite explorar as poses possíveis dos compostos selecionados para prospecção nos receptores associados à pele e atribuir uma pontuação a cada pose. Essa triagem fornece critérios para priorizar moléculas e orientar análises estruturais, contribuindo para investigar a diversidade das substâncias produzidas por microalgas \parencite{khemiri2023,trott2010,warren2006}.

Vina e AD4 calculam a FEB por funções distintas. Vina combina termos de ajuste estérico, contatos hidrofóbicos, ligações de hidrogênio e flexibilidade do ligante; AD4 usa sua própria parametrização para interações de van der Waals, ligações de hidrogênio, eletrostática, dessolvatação e torção. Nos alvos metaloproteicos, AD4Zn acrescenta parâmetros para representar a coordenação com zinco. Embora os valores sejam expressos em kcal/mol, cada função produz sua própria escala: uma FEB mais negativa em AD4 do que em Vina não demonstra maior afinidade. Por isso, comparam-se ligantes apenas no mesmo receptor e sob a mesma função; a FEB ordena poses segundo o modelo empregado, sem medir afinidade experimental ou inibição enzimática \parencite{morris2009autodock,santos2014zn,trott2010,warren2006}.

A visualização molecular no PyMOL e os contatos identificados pelo PLIP ajudam a examinar as poses e os resíduos envolvidos. Quando disponíveis, os ligantes coocristalizados oferecem uma referência estrutural para avaliar o posicionamento no sítio \parencite{pymol,salentin2015plip,cole2005}. A leitura conjunta dessas ferramentas permite comparar os compostos selecionados para prospecção com os ligantes comerciais e priorizar os perfis de interação de maior interesse para estudos posteriores.
